% Pursuing Stacks (cotes 134-2 to 134-8): archive cote 134-6, site page 51 (PDF page 52), typescript page 369.
% typescript. Draft compared: ps5.tex:3898-4016.
% AI-transcribed by Claude Opus 5.5 (claude-opus-5-5, adaptive thinking, effort high), Message Batches API, batch ps3 (msgbatch_01XdFmuVVBAcgc7EFxdtCqsj), 2026-11-12; not yet checked by eye.
\scanpage{134-6/51}
\mspage{369}
(10) \qquad $\mathrm{R}\Gamma_B(K^\bullet_B) \hand{$\simeq$} \mathrm{Hom}^{\bullet\bullet}(p^{ab}_{B!}(L^B_\bullet), K^\bullet)$ \quad ,
where this time the Hom's in the \xout{left} right hand side of (10) are taken
in (Ab), not in $B^\wedge_{ab}$ .

This formula very strongly suggests to view the chain complex of
abelian groups
(11) \qquad $p^{ab}_{B!}(L^B_\bullet)$ \quad ,
which is in fact a complex of projective (hence free) abelian groups
defined up to chain homotopy, as embodying the global homology structure
of B (or of the corresponding topos), more accurately still, as embodying
the homology structure of the corresponding homotopy type. It is easily
seen that the corresponding object of D.(Ab) depends covariantly on B
when B varies in the category (Cat), so that we get a functor
\[ (\mathrm{Cat}) \hand{$\longrightarrow$} \mathrm{D.(Ab)} \overset{\mathrm{def}}{=} (\mathrm{Hotab}) \quad , \]
which in view of (10) (an isomorphism functorial not only \xout{inxK} with respect
to $K^\bullet$ , but equally with respect to B ) factors through the localization
(Hot) of (Cat), thus yielding a canonical functor
(12) \qquad $(\mathrm{Hot}) \hand{$\longrightarrow$} (\mathrm{H}\overset{\hand{$\smile$}}{\underset{\hand{$\frown$}}{\mathrm{o}}}\mathrm{tab})$ \quad ,
which deserves to be called the \uline{abelianization functor}\hand{,} from homotopy
types to ``abelian homotopy types'' . This cannot be of course anything else
(up to canonical isomorphism) but the functor (1) of section 92 (p.321),
but obtained here in a\xout{n} wholly ``intrinsic'' way, without having to pass
through the particular properties of a particular test category such as
\hand{$\Delta$} or one of it's twins. One possible way to check this identity would
be by proving that \deleted{the same} \inserted{an} isomorphism (10) is valid when replacing
(for a given B in (Cat) and $K^\bullet$ in $D^+$(Ab)) the chain complex (11) by the
corresponding one deduced from the map (1) defined p.321 (via the diagram
(3) on p.322), and checking moreover that an object \inserted{$\ell_\bullet$} of $D^-$(Ab) is known
up to canonical isomorphism, when we know the corresponding functor
(13) \qquad $K^\bullet \hand{$\longmapsto$} \mathrm{Hom}_{D(Ab)}(\hand{$\ell_\bullet$} , K^\bullet)$
on $D^+$(Ab). Presumably this latter statement holds when replacing \xout{Ab} (Ab)
by any abelian category, but I confess I did\unsure{'}nt sit down to check it, nor
do I remember having seen it stated somewhere - as I do'nt remember either
having seen anywhere a comprehensive treatment about the relationship bet-%
ween homology and cohomology. So maybe my present reflections do fill a
gap, or at any rate give some indications as to how to fill it...

\margin{51}
